\section{我怎么理解这道题，以及先砍掉了什么}

题目给的场景是宠物对战游戏：玩家养一批精灵，打 PVP，配阵容，日常培养。要交付的是一个 LLM 教练，具备军师、陪练、老师三种角色感，至少完整闭环其中一项，其余可以只做设计。参考方向列了 Agentic RL、RAG、ReAct、Memory Stream / Reflection，产品参考是和平精英的"小田"和王者荣耀的"灵宝"。

拿到题目我先做的不是写代码，而是划边界。三件事我判断为本次不做，理由写在下面，因为不做的决定和做的决定一样需要交代。

\textbf{不做三渲二的美术和完整游戏内容。}题目要的是教练，游戏是产生局面的环境。我把宠物控制在 12 只、5 个关卡，够产生"抢先收尾还是留能量""换宠承伤还是原地治疗"这类真实取舍就停手。继续堆宠物不会让教练更聪明，只会让平衡工作吃掉全部时间。

\textbf{不做 LLM 权重训练。}本机是 arm64、没有可用的 CUDA 后端，也没有云上 GPU 预算。我做了两个能在这台机器上真跑完的训练实验（后面第 4 节和第 6 节会讲它们各自是什么、不是什么），但没有把它们包装成"训练了模型"。

\textbf{不做联网 PVP。}本地对战我做成了同一台设备上分屏轮流操作。它离真实竞技还有距离，这一点我在第 5 节里明说了，没有含糊过去。

三个角色里，军师和老师我做到了完整闭环，陪练只做到设计。这个划分是刻意的，也符合题目"其它可仅做设计"的说法。但我在第 3 节会把话说清楚，避免读的人以为三样都做完了。

\section{怎么看这个东西}

本地跑起来：

\begin{verbatim}
cd pet-coach-game
npm start          # Node.js 20+，无第三方依赖
\end{verbatim}

浏览器打开 \texttt{http://127.0.0.1:8765/}。想看模型能力需要先在 \texttt{/connect.html} 里录入 DeepSeek 密钥，密钥只留在后端进程内存里；不录也能玩，规则建议、复盘、培养建议和小测都不依赖模型。

进页面是营地，顶部有两个入口：\textbf{训练}和\textbf{对局}。训练按关卡打固定对手，教练全程给建议；对局里对手从全部 12 只里自动配队、等级按你的队伍适配，教练照常工作。对局还有第二种对手：真人同机，两个人分屏各选各的，两边各有一条教练。

想快速看到教练在做什么，有两处最直接：一是战斗页底部左右两侧的教练条，各自只分析自己那一侧的局面；二是打完一局后自动出现的整局复盘，它会指出具体是哪一回合、当时有哪些别的选择。

\section{三个角色各自做到哪一步}

\subsection{军师：完整闭环}

军师是这次的重心。它的工作方式是把"算"和"说"分开：所有伤害、行动顺序、合法性由游戏引擎枚举，模型不参与计算，只负责把比较结果讲成玩家能理解的取舍。

具体做法是，玩家面对一个局面时，引擎枚举双方所有合法行动，对每一种组合跑一遍共享结算器，得到平均收益和最坏分支。同时顺序（双方同速）会分别按两种可能各算一次，这样搜索不会偷偷利用随机数。模型拿到的证据里包含这些数字、每张战术卡的适用条件和反例，它要做的只是解释。

举例说明这个分工为什么必要。同一个残血目标，火花直接伤害 52，追猎 75，两个都够击倒，都耗 2 豆。如果只看面板数字，很自然会告诉玩家"选追猎"。但这个判断是错的：对方可能吃药、可能换宠、可能防御，而追猎的额外伤害来自"目标已灼烧"这个前提，前提一变结论就变。这类判断交给模型去猜是猜不准的，交给枚举器是确定的。

\subsection{老师：完整闭环}

复盘分成整局和单回合两层。整局先给行动分布和回合数，再挑几个造成减员或血量大幅变化的回合；单回合用回合前的局面做条件比较，比如"当时目标 78 血，火花不防御打 24、防御打 8，单次攻击不足以收尾"。

关键的一条是：复盘一律使用回合开始时的公开状态，不使用对手实际打出的那一招。理由很直白，如果拿事后结果当依据，任何一次失败都能被解释成"你当时选错了"，而这既不公平也没法教人。这条规则有测试保证，改对手的实际行动不会改变复盘结论。

练习跟着复盘走。玩家在某个回合做了选择，教练按这个选择对应的知识点出一道题（灼烧追击、防御节奏、换宠承伤、先制与速度四类）。出题、等待作答、判题、取消这四个状态由程序控制，模型不能提前报答案，也不能把题目改写成别的。这道题和玩家刚才的真实选择挂钩，不是随机题库。

\subsection{陪练：只做了设计}

这是我这次最薄的一块，我把现状说清楚。

已经实现的只有：玩家表达挫败时的克制回应、追问时接住上文、以及基于真实对战历史的开场白（"记得上一场是在青芽草地"）。代码量很小。

没有实现的是"有情绪"这件事本身。我写了设计文档，里面做了一套情绪状态和对应的语气档位，但代码里没有。所以如果问"陪练有没有情绪"，答案是：设计上有，实现上没有。

为什么没做？两个原因。一是题目允许非主角色只做设计，我把时间投给了军师和老师；二是我对"给游戏助手加情绪"这件事本身有保留。项目调研过的灵宝资料里有一条我认同的判断：情绪陪伴很容易变成油腻，玩家不打开聊天不代表需要被安慰。所以我的设计文档里给它加了比一般助手更严的约束——允许沉默、不主动提问、不空泛安慰、不因为一次失败就弹话。
